\centerline{\bf \Huge Антигеомный разнобой}

\problem{1}
Рассмотрим график приведенного квадратного трехчлена с положительным
дискриминантом. Проведем окружность через три точки пересечения этого
графика с осями координат. Докажите, что все такие окружности проходят
через фиксированную точку на плоскости.

\problem{2}
 языке племени УЫ всего две буквы: У и Ы. Словом считается любая
последовательность из $2 n$ букв У и $2 n$ букв Ы (число $n$ дано и
фиксировано). Язьковеды называют слова похожими, если одно можно
получить из другого одной перестановкой двух соседних букв У и Ы. Какое
наибольшее количество слов можно выписать на доску так, чтобы любые два
из выписанных слов не были похожи?

\problem{3}
Пусть $a$ и $b$ — натуральные числа, такие, что число $a+b^3$ делится на
число $a^2+3ab+3b^2-1$. Докажите, что число $a^2+3ab+3b^2-1$ делится на
точный куб натурального числа, большего, чем 1.

\problem{4}
Дано натуральное $n$. В ряд записаны $n$ целых чисел. Два человека
играют в игру. За один шаг первый указывает на несколько из них,
записанных подряд, а второй либо увеличивает каждое из указанных чисел
на 1, либо уменьшает каждое из них на 1. Найдите наибольшее $k$, такое,
для которого первый всегда за несколько шагов сможет добиться, чтобы
хотя бы $k$ чисел стали делиться на 3.

\problem{5}
Найти наибольшее натуральное $n$, при котором сумма
$[\sqrt{1}]+[\sqrt{2}]+\ldots+[\sqrt{n}]$ является простым числом.

\problem{6}
Дана доска $2m\times 2n$, раскрашенная в шахматном порядке. Сколькими
способами можно поставить на белые клетки $mn$ шашек так, чтобы никакие
две шашки не стояли бы на одной клетке и чтобы никакие две шашки не
располагались бы в клетках, соседних по углу?

\problem{7}
Алиса и Боб играют в игру. На плоскости отмечены $n$ точек общего
положения, где $n$ — нечетное натуральное число; $n$ — произвольное
натуральное число, большее 1. Алиса и Боб по очереди выбирают пару точек
и соединяют их отрезком. Разрешается соединять только точки, между
которыми еще не проведен отрезок. Проигрывает тот, после чьего хода
образуется цикл нечетной длины. Начинает Алиса. Кто выигрывает при
правильной игре?


\newpage

\centerline{}

\problem{8}
Дан многочлен $P(x)=ax^3+bx^2+cx+d$ с целыми коэффициентами, такими, что
$a\neq 0$. Известно, что существует бесконечно много пар целых чисел
$(x,y)$, для которых справедливо равенство $xP(x)=yP(y)$. Докажите, что
у многочлена $P(x)$ есть целый корень.